\documentclass[a4paper,11pt]{article}
\usepackage[a4paper, top=2.5cm, bottom=2.5cm, left=2cm, right=2cm]{geometry}
\usepackage{enumitem}
\setlength\parindent{0pt}
\usepackage{fancyhdr}


\pagestyle{fancy}
\fancyhf{}
\fancyfoot[L]{TN High-Tech Physics}
\fancyfoot[C]{\today}
\fancyfoot[R]{2025-2026}

\renewcommand{\footrulewidth}{0.4pt}
\renewcommand{\headrulewidth}{0pt}

\begin{document}

\section*{Opdrachtbeschrijving Memo 1 -- Technische Stage}
\subsection*{Introductie van de technology \textbf{(20\%)}}
Geef een korte inleiding over de technology waaraan jouw team werkt. Behandel ten minste:
\begin{itemize}[leftmargin=1.5cm, partopsep=1pt, itemsep=2pt, parsep=0pt]
    \item het onderzoeksveld waarin de technologie wordt toegepast;
    \item het werkingsprincipe van de technologie;
    \item de (eventuele) problemen die de technologie kan oplossen;
    \item verdere ontwikkelingen en uitdagingen.
\end{itemize}

\textbf{Criterium:} Een medestudent die niet bekend is met het project moet het werkingsprincipe en de bredere context van de technology kunnen begrijpen.

\subsection*{Doel van de technische stage \textbf{(20\%)}}
Beschrijf het doel van jouw stage. Beantwoord hierbij:
\begin{itemize}[leftmargin=1.5cm, partopsep=1pt, itemsep=2pt, parsep=0pt]
    \item Wat wil je aan het einde van de stage hebben bereikt?
    \item Hoe is het doel afgebakend? Zijn er subdoelen? Bespreek dit met je begeleider.
    \item Wat is jouw specifieke rol binnen het behalen van dit doel?
\end{itemize}

\textbf{Criterium:} De lezer moet begrijpen wat het doel van de stage is, inclusief de afbakening en jouw rol daarin.

\subsection*{Planning, overdracht en labjournaal \textbf{(20\%)}}
Beschrijf in deze sectie de planning (per student) voor het behalen van het stagedoel. Begin met een overzicht van de activiteiten die je gaat uitvoeren, inclusief de volgorde en de geschatte tijdsduur. Denk bijvoorbeeld aan:
\begin{itemize}[leftmargin=1.5cm, partopsep=1pt, itemsep=2pt, parsep=0pt]
    \item kalibratiemetingen;
    \item ontwerp en bouw van een opstelling;
    \item schrijven en testen van software.
\end{itemize}
Geef aan welke stappen nog moeten worden uitgezocht of voorbereid, zodat de lezer kan inschatten hoe compleet het plan is. Leg daarnaast uit hoe jouw werkzaamheden na de stage bruikbaar zijn voor de vakgroep:
\begin{itemize}[leftmargin=1.5cm, partopsep=1pt, itemsep=2pt, parsep=0pt]
    \item structuur van data en bestanden;
    \item gebruik van versiebeheer (bijv. GitHub);
    \item centrale opslag (bijv. Google Drive, Overleaf, OneNote,labjournaal);
    \item overzicht van afgeronde taken zodat toekomstige studenten of onderzoekers hierop kunnen voortbouwen.
\end{itemize}

Denk na over het gebruik van een (digitaal) labjournaal:
\begin{itemize}[leftmargin=1.5cm, partopsep=1pt, itemsep=2pt, parsep=0pt]
    \item notatie van meetcondities, parameters en resultaten;
    \item koppeling tussen labjournaal en dataopslag voor reproduceerbaarheid.
\end{itemize}

\textbf{Criterium:} De lezer moet kunnen volgen hoe de planning van de stage verloopt, hoe de werkzaamheden overdraagbaar zijn en hoe experimentel handelingen genoteerd worden in het labjournaal.

\subsection*{Opstelling en eerste resultaten \textbf{(20\%)}}
Beschrijf de opstelling of het experiment en eventuele eerste resultaten. Resultaten mogen ruim worden opgevat, bijvoorbeeld:
\begin{itemize}[leftmargin=1.5cm, partopsep=1pt, itemsep=2pt, parsep=0pt]
    \item schets van de opstelling;
    \item gebruikte apparatuur;
    \item eerste observaties, grafieken of tabellen;
    \item tests van software of kalibratiemetingen.
\end{itemize}

\textbf{Criterium:} De opstelling en eerste resultaten moeten volledig en duidelijk beschreven zijn, zodat een medestudent het experiment kan begrijpen en reproduceren.

\subsection*{Algemene criteria \textbf{(20\%)}}
\begin{itemize}[leftmargin=1.5cm, partopsep=1pt, itemsep=2pt, parsep=0pt]
    \item Lengte: 2-3, maximaal 4 A4;
    \item Grammaticaal en taalkundig correct;
    \item Nette en overzichtelijke lay-out;
    \item Op tijd ingeleverd;
    \item Geen plagiaat;
    \item Correcte referenties \& gebruik AI (inclusief prompts).
\end{itemize}

\end{document}
